\documentclass[12pt]{article}
\usepackage[utf8]{inputenc}
\usepackage[T1]{fontenc}
\usepackage[margin=1in]{geometry}
\usepackage{setspace}
\usepackage{xcolor}
\usepackage[colorlinks=true, urlcolor=blue!60!black, linkcolor=blue!60!black]{hyperref}
\setlength{\parindent}{0em}
\setlength{\parskip}{0.6em}
\renewcommand{\baselinestretch}{1.0}
\begin{document}
\begin{center}
{\LARGE \textit{Meyer v. Nebraska}, 262 U.S. 390 (1923)\par}
\vspace{0.5em}
{\large Case Brief --- Two-Page Edition\par}
\end{center}

\section{Memory Jogger}
Nebraska's ban on teaching modern foreign languages to young children violated Fourteenth Amendment ``liberty,'' which protects teachers, pupils, and parents' right to direct their children's education (\href{https://www.courtlistener.com/opinion/100233/meyer-v-nebraska/}{CourtListener}).

\section{Detailed Case Facts}
After World War I, Nebraska's Siman Act forbade teaching any subject in a language other than English, or teaching modern foreign languages to children who had not passed the eighth grade. Robert T. Meyer, a teacher at a Lutheran parochial school, was charged for teaching a ten-year-old to read German using Biblical stories.

\section{Procedural History}
Meyer was convicted and fined \$25. The Nebraska Supreme Court affirmed, upholding the act as a police-power measure to make English the ``mother tongue'' of all children.

\section{Judicial Votes}
Reversed 7--2. McReynolds wrote for Taft, McKenna, Van Devanter, Brandeis, Butler, and Sanford; Holmes, joined by Sutherland, dissented in \textit{Bartels v. Iowa}.

\section{Holding}
The act arbitrarily infringed the liberty protected by the Due Process Clause, resting on the teacher's right to teach and the parents' right to engage him. The broad catalog of liberties, and the reserved power to require English instruction and set public curricula, were dicta.

\section{Analysis of Opinions}
\textbf{Majority:} Liberty reaches knowledge, marriage, child-rearing, and worship. ``[M]ere knowledge of the German language cannot reasonably be regarded as harmful,'' and the wish ``to foster a homogeneous people with American ideals'' could not justify forbidden means; Plato's and Sparta's state-reared children are alien to American institutions. \textbf{Holmes:} a common tongue is a lawful end, and the question was one ``upon which men reasonably might differ.''

\section{Examples: Future Applications}
\textbf{Same-Side:} compulsory public schooling that outlaws private schools falls (\textit{Pierce v. Society of Sisters}). \textbf{Opposite-Side:} requiring every school to teach English is valid. \textbf{Fence-Sitter:} testing requirements for homeschoolers or limits on curriculum opt-outs.

\section{Critique}
Originalists see a \textit{Lochner}-era invention that seeded modern substantive due process. Progressives praise its defense of pluralism, while Barbara Bennett Woodhouse argues it treated ``the child as property'' of parents.

\section{Key Quotations}
\begin{quote}
``Without doubt, it denotes not merely freedom from bodily restraint, but also the right of the individual to contract, to engage in any of the common occupations of life, to acquire useful knowledge, to marry, establish a home and bring up children, to worship God according to the dictates of his own conscience, and generally to enjoy those privileges long recognized at common law as essential to the orderly pursuit of happiness by free men.''
\end{quote}
\noindent --- Justice McReynolds, Opinion of the Court, 262 U.S. at 399.

\end{document}
